\section*{अध्याय \ref{ch:g7:sediments-and-fossils} --- अवसाद और जीवाश्म}

\begin{solution}{exo:g7:sediments-and-fossils:1}
(क) \omterm{def:g7:sediments-and-fossils:fossil}{साँचा}; (ख) \omterm{def:g7:sediments-and-fossils:fossil}{चिह्न जीवाश्म}; (ग) \omterm{def:g7:sediments-and-fossils:sediment}{अवसाद}; (घ) \omterm{def:g7:sediments-and-fossils:fossil}{ढलाई}; (ङ) \omterm{def:g7:sediments-and-fossils:layer}{स्तर}; (च)
\omterm{def:g7:sediments-and-fossils:fossil}{जीवाश्म}।
\end{solution}

\begin{solution}{exo:g7:sediments-and-fossils:2}
(ख) जानवर मरता है; (घ) \omterm{def:g3:stones-pebbles-sand:sand}{कीचड़} सीपी को जल्दी ढक लेता है; (ग) नई परतें जमा होती हैं
और सीपी पत्थर बन जाती है; (क) \omterm{def:g5:water-shapes-the-land:erosion}{अपरदन} उसके ऊपर की परतों को घिस
देता है।
\end{solution}

\begin{solution}{exo:g7:sediments-and-fossils:3}
(क) एक गरम, उथला, साफ़ \omterm{def:g2:water-on-earth:ocean}{समुद्र}। (ख) उथला पानी, \omterm{def:g1:landscapes-around-us:water}{झील}, \omterm{def:g1:landscapes-around-us:water}{नदी} या \omterm{def:g2:water-on-earth:ocean}{समुद्र} के किनारे।
(ग) एक दलदली जंगल या दलदल। (घ) ज़मीन, जिस पर पौधे उगते थे (और व्याध पतंगे उड़ते
थे), पानी के पास जहाँ \omterm{def:g3:stones-pebbles-sand:sand}{कीचड़}
बैठता था।
\end{solution}

\begin{solution}{exo:g7:sediments-and-fossils:4}
(क) ग़लत: ज़्यादातर जीव कोई \omterm{def:g7:sediments-and-fossils:fossil}{जीवाश्म} नहीं छोड़ते; दबने से पहले वे सड़ जाते हैं या
घुल जाते हैं। (ख) सही: यह एक \omterm{def:g7:sediments-and-fossils:fossil}{चिह्न जीवाश्म} है। (ग) ग़लत: सबसे बड़े और भारी कण
पहले बैठते हैं, महीन \omterm{def:g3:stones-pebbles-sand:sand}{कीचड़} सबसे आख़िर में। (घ) सही।
\end{solution}

\begin{solution}{exo:g7:sediments-and-fossils:5}
(क) बलुआ पत्थर \qty{120}{m}, चूना पत्थर \qty{180}{m}, शेल \qty{100}{m}।
(ख) $120 + 180 + 100 = \qty{400}{m}$। (ग) $180 \div 400 = 0.45$:
\textbf{45~\%}। (घ) उसका ऊपरी हिस्सा \omterm{def:g5:water-shapes-the-land:erosion}{अपरदन} से घिसकर हट गया है: जो बचा है वह पूरी
परत नहीं है, और कितना बचा है यह जगह-जगह बदलता है।
\end{solution}

\begin{solution}{exo:g7:sediments-and-fossils:6}
(क) नीचे से ऊपर: \omterm{def:g3:stones-pebbles-sand:pebble}{बजरी}, फिर \omterm{def:g3:stones-pebbles-sand:sand}{रेत}, फिर सबसे ऊपर \omterm{def:g3:stones-pebbles-sand:sand}{कीचड़}, साफ़ पानी के नीचे। (ख) \omterm{def:g3:stones-pebbles-sand:sand}{कीचड़}
के नन्हे कण बहुत धीरे-धीरे डूबते हैं: एक घंटे बाद भी वे गिर रहे होते हैं। (ग) पहली
परतों के ऊपर परतों का एक दूसरा समूह: फिर से \omterm{def:g3:stones-pebbles-sand:sand}{रेत}, फिर \omterm{def:g3:stones-pebbles-sand:sand}{कीचड़}। (घ) बिल्कुल नीचे की
\omterm{def:g3:stones-pebbles-sand:pebble}{बजरी}।
\end{solution}

\begin{solution}{exo:g7:sediments-and-fossils:7}
सबसे अच्छी: (ख), एक सख़्त खोल, जो शांत पानी में महीन \omterm{def:g3:stones-pebbles-sand:sand}{कीचड़} के नीचे जल्दी दब गया।
फिर (घ): सीपी सख़्त है और दबी है, पर \omterm{def:g5:water-shapes-the-land:flood}{बाढ़} उसे हिला और तोड़ भी सकती है। सबसे कम:
(क), बिना सख़्त हिस्सों वाला एक नरम शरीर, \omterm{def:g3:stones-pebbles-sand:rock}{चट्टानों} पर जहाँ उसे ढकने के लिए कोई
\omterm{def:g3:stones-pebbles-sand:sand}{कीचड़} नहीं बैठता। (ग) भी ज़्यादा बेहतर नहीं: हड्डी खुली पड़ी रहती है और टूटकर घिस
जाती है।
\end{solution}

\begin{solution}{exo:g7:sediments-and-fossils:8}
(क) एक खोखली जगह जिसकी दीवारों पर सीपी का आकार छपा है: एक \omterm{def:g7:sediments-and-fossils:fossil}{साँचा}। (ख) एक \omterm{def:g7:sediments-and-fossils:fossil}{ढलाई},
सीपी का पत्थर का प्रतिरूप। (ग) नहीं: सीपी घुल चुकी है; \omterm{def:g7:sediments-and-fossils:fossil}{ढलाई} नए \omterm{def:g6:rocks-and-minerals:mineral}{खनिजों} से बनी है।
(घ) पानी \omterm{def:g3:stones-pebbles-sand:sand}{रेत} के बड़े कणों के बीच से आसानी से गुज़रता है, इसलिए वह सीपी तक पहुँचकर
उसे घोल सकता है; महीन \omterm{def:g3:stones-pebbles-sand:sand}{कीचड़} में वह कहीं कम आसानी से
गुज़रता है।
\end{solution}

\begin{solution}{exo:g7:sediments-and-fossils:9}
(क) $1.6 \times 1000 = \qty{1600}{m}$। (ख) $1600 \div 80 = 20$ परतें।
(ग) $160 \div 1600 = 0.1$: \textbf{10~\%}।
\end{solution}

\begin{solution}{exo:g7:sediments-and-fossils:10}
(क) नहीं: पैर केवल नरम ज़मीन, गीली \omterm{def:g3:stones-pebbles-sand:sand}{रेत} या \omterm{def:g3:stones-pebbles-sand:sand}{कीचड़} में निशान छोड़ता है; ज़मीन बाद में
\omterm{def:g3:stones-pebbles-sand:rock}{चट्टान} बनी। (ख) अगली लहर या \omterm{def:g5:water-shapes-the-land:flood}{बाढ़} उन्हें धो देती है। (ग) नए \omterm{def:g7:sediments-and-fossils:sediment}{अवसाद} को उन्हें मिटने से
पहले जल्दी ढकना और भरना था। (घ) \omterm{def:g7:sediments-and-fossils:fossil}{चिह्न जीवाश्म}: यह दिखाता है कि जानवर ने क्या किया,
उसके शरीर का कोई हिस्सा नहीं।
\end{solution}

\begin{solution}{exo:g7:sediments-and-fossils:11}
(क) $20 + 5 = \qty{25}{mm}$। (ख) $1000 \div 25 = 40$ जोड़ियाँ: कोर
\textbf{40 \omterm{def:g5:water-shapes-the-land:flood}{बाढ़ें}} दर्ज करता है। (ग) \omterm{def:g3:stones-pebbles-sand:sand}{रेत} भारी है: \omterm{def:g1:landscapes-around-us:water}{नदी} के \omterm{def:g1:landscapes-around-us:water}{झील} में आने के जल्दी बाद,
\omterm{def:g5:water-shapes-the-land:source-mouth}{मुहाने} के पास, वह बैठ जाती है; उसका थोड़ा ही हिस्सा दूर तक जाता है। (घ) दूर, पानी
में केवल महीन \omterm{def:g3:stones-pebbles-sand:sand}{कीचड़} बचता है: वह वहाँ धीरे-धीरे, पानी शांत होने पर,
बैठता है।
\end{solution}

\begin{solution}{exo:g7:sediments-and-fossils:12}
(क) सीपियाँ सख़्त होती हैं: वे दबने तक काफ़ी समय टिकी रहती हैं, और सीपी वाले बहुत-से
जानवर \omterm{def:g2:water-on-earth:ocean}{समुद्र} की तली पर रहते थे, जहाँ उन पर \omterm{def:g3:stones-pebbles-sand:sand}{कीचड़} बैठता है। (ख) नरम शरीर दबने से
पहले सड़ जाते हैं। (ग) नहीं: नरम जानवर लगभग कभी \omterm{def:g7:sediments-and-fossils:fossil}{जीवाश्म} नहीं छोड़ते, इसलिए
\omterm{def:g7:sediments-and-fossils:fossil}{जीवाश्मों} से बनी सूची अधूरी है; वह ज़्यादातर सीपियों या हड्डियों वाले जानवर
दिखाती है।
\end{solution}

\begin{solution}{pb:g7:sediments-and-fossils:1}
\textbf{1.} बलुआ पत्थर, चूना पत्थर, शेल, खड़िया (पहली स्तंभ में सबसे
नीचे)।

\textbf{2.} \qty{120}{m}, \qty{180}{m} और \qty{100}{m}; कुल
\qty{400}{m}।

\textbf{3.} $180 \div 400 = 0.45$: \textbf{45~\%}।

\textbf{4.} उसका ऊपरी हिस्सा घिसकर हट गया है: परत का ऊपरी भाग ग़ायब है, और जो अब
वहाँ है ही नहीं उसे कोई माप नहीं सकता।

\textbf{5.} \omterm{def:g1:landscapes-around-us:water}{झील}, \omterm{def:g1:landscapes-around-us:water}{नदी} या \omterm{def:g2:water-on-earth:ocean}{समुद्र} के किनारे उथले पानी में: तरंग-चिह्न उथले पानी में \omterm{def:g3:stones-pebbles-sand:sand}{रेत}
पर बनते हैं, और पैरों के निशान नरम, गीली \omterm{def:g3:stones-pebbles-sand:sand}{रेत} या \omterm{def:g3:stones-pebbles-sand:sand}{कीचड़} में, अक्सर किसी \omterm{def:g1:landscapes-around-us:water}{तट} पर, छूटते
हैं।

\textbf{6.} एक गरम, उथले, साफ़ \omterm{def:g2:water-on-earth:ocean}{समुद्र} में: प्रवाल।

\textbf{7.} शांत: महीन \omterm{def:g3:stones-pebbles-sand:sand}{कीचड़} केवल वहाँ बैठता है जहाँ पानी लगभग स्थिर है।

\textbf{8.} पहले वह जगह एक \omterm{def:g1:landscapes-around-us:water}{तट} थी, उथले पानी और नरम \omterm{def:g3:stones-pebbles-sand:sand}{रेत} वाली, जहाँ जानवर चलते थे।
फिर \omterm{def:g2:water-on-earth:ocean}{समुद्र} ने उसे ढक लिया: एक गरम, उथला \omterm{def:g2:water-on-earth:ocean}{समुद्र} जहाँ प्रवाल बढ़े और सीपियाँ जमा
हुईं। बाद में पानी और शांत हो गया, शायद और गहरा, और केवल महीन \omterm{def:g3:stones-pebbles-sand:sand}{कीचड़} बैठा,
अमोनाइट की सीपियों के साथ, जो स्क्विड और ऑक्टोपस के संबंधी
जानवर थे।

\textbf{9.} $180 \times 100 = 18\,000$~सेमी।

\textbf{10.} $120~\text{m} = 12\,000$~सेमी, और $12\,000 \times 20 =
240\,000$~वर्ष।

\textbf{11.} $1~\text{m} = 100$~सेमी, और $100 \times 100 = 10\,000$~वर्ष।

\textbf{12.} $180 \times 10\,000 = 1\,800\,000$: लगभग
\textbf{1\,800\,000 वर्ष}, लगभग दो मिलियन वर्ष। यह बलुआ पत्थर से
$1\,800\,000 \div 240\,000 = 7.5$ गुना ज़्यादा है।
\end{solution}
